% ============================================================
%  Lost in Translation, Lost in Safety
%  ACLing 2026 — Procedia Computer Science
%  Stage 6 draft — simulation results (camera-ready: real API)
%  Double-blind submission (author info withheld)
% ============================================================
\documentclass[3p,times,procedia]{elsarticle}
\flushbottom

\usepackage{ecrc}
\usepackage[bookmarks=false]{hyperref}
  \hypersetup{colorlinks,linkcolor=blue,citecolor=blue,urlcolor=blue}
\usepackage{amssymb}
\usepackage{booktabs}
\usepackage{xcolor}
\usepackage{multirow}
\usepackage{graphicx}
\usepackage{amsmath}
\usepackage[figuresright]{rotating}
\usepackage{scalefnt}

% Define \doi for Vancouver bibliography style
\providecommand{\doi}[1]{\href{https://doi.org/#1}{\nolinkurl{doi:#1}}}

% ── Locked result macros (Stage 5) ──────────────────────────
% macro_defs.tex — auto-generated by generate_tables.py
% All simulation=True. Do NOT edit manually.
% Include with: % macro_defs.tex — auto-generated by generate_tables.py
% All simulation=True. Do NOT edit manually.
% Include with: % macro_defs.tex — auto-generated by generate_tables.py
% All simulation=True. Do NOT edit manually.
% Include with: \input{macro_defs}

\newcommand{\GptFourOEnAcc}{0.863} % GPT-4o EN accuracy
\newcommand{\MistralEnAcc}{0.627} % Mistral-7B EN accuracy
\newcommand{\GptFourOEsEnAcc}{0.7638} % GPT-4o ES-EN mean accuracy
\newcommand{\GptFourOArEnAcc}{0.7486} % GPT-4o AR-EN mean accuracy
\newcommand{\HThreeRho}{-1.000} % Spearman rho CS dose-response
\newcommand{\HThreeP}{0.0000} % Spearman p-value
\newcommand{\GptMistralCsGap}{0.301} % GPT-4o minus Mistral-7B ES-EN accuracy gap
\newcommand{\QwenArDrop}{0.122} % Qwen2-72B AR-EN accuracy drop
\newcommand{\LlamaArDrop}{0.172} % LLaMA-3.1-70B AR-EN accuracy drop
\newcommand{\NItemsTotal}{5{,}200} % Total benchmark items
\newcommand{\NItemsBase}{400} % EN base items
\newcommand{\NPartitions}{13} % Benchmark partitions
\newcommand{\NModels}{7} % Evaluated LLMs
\newcommand{\NRuns}{3} % Evaluation runs per model/partition

\newcommand{\GptFourOEnAcc}{0.863} % GPT-4o EN accuracy
\newcommand{\MistralEnAcc}{0.627} % Mistral-7B EN accuracy
\newcommand{\GptFourOEsEnAcc}{0.7638} % GPT-4o ES-EN mean accuracy
\newcommand{\GptFourOArEnAcc}{0.7486} % GPT-4o AR-EN mean accuracy
\newcommand{\HThreeRho}{-1.000} % Spearman rho CS dose-response
\newcommand{\HThreeP}{0.0000} % Spearman p-value
\newcommand{\GptMistralCsGap}{0.301} % GPT-4o minus Mistral-7B ES-EN accuracy gap
\newcommand{\QwenArDrop}{0.122} % Qwen2-72B AR-EN accuracy drop
\newcommand{\LlamaArDrop}{0.172} % LLaMA-3.1-70B AR-EN accuracy drop
\newcommand{\NItemsTotal}{5{,}200} % Total benchmark items
\newcommand{\NItemsBase}{400} % EN base items
\newcommand{\NPartitions}{13} % Benchmark partitions
\newcommand{\NModels}{7} % Evaluated LLMs
\newcommand{\NRuns}{3} % Evaluation runs per model/partition

\newcommand{\GptFourOEnAcc}{0.863} % GPT-4o EN accuracy
\newcommand{\MistralEnAcc}{0.627} % Mistral-7B EN accuracy
\newcommand{\GptFourOEsEnAcc}{0.7638} % GPT-4o ES-EN mean accuracy
\newcommand{\GptFourOArEnAcc}{0.7486} % GPT-4o AR-EN mean accuracy
\newcommand{\HThreeRho}{-1.000} % Spearman rho CS dose-response
\newcommand{\HThreeP}{0.0000} % Spearman p-value
\newcommand{\GptMistralCsGap}{0.301} % GPT-4o minus Mistral-7B ES-EN accuracy gap
\newcommand{\QwenArDrop}{0.122} % Qwen2-72B AR-EN accuracy drop
\newcommand{\LlamaArDrop}{0.172} % LLaMA-3.1-70B AR-EN accuracy drop
\newcommand{\NItemsTotal}{5{,}200} % Total benchmark items
\newcommand{\NItemsBase}{400} % EN base items
\newcommand{\NPartitions}{13} % Benchmark partitions
\newcommand{\NModels}{7} % Evaluated LLMs
\newcommand{\NRuns}{3} % Evaluation runs per model/partition

% ── Procedia header ─────────────────────────────────────────
\volume{00}
\firstpage{1}
\journalname{Procedia Computer Science}
\runauth{Anonymous et al.}
\jid{procs}

% ── Simulation disclaimer shortcut ──────────────────────────
\newcommand{\simflag}{\textsuperscript{\textdagger}}
\newcommand{\simtext}{%
  \textsuperscript{\textdagger}All numerical results derive from
  simulation mode (sandbox constraints; no live API keys). Camera-ready
  version will replace these with real API evaluations.}

\begin{document}
\begin{frontmatter}

\dochead{ACLing 2026: 8th International Conference on AI in Computational Linguistics}

\title{Lost in Translation, Lost in Safety: Benchmarking Multilingual LLMs
       for Consumer Medication Risk Understanding}

% ── DOUBLE-BLIND: author information withheld ───────────────
\author[anon]{Anonymous Author(s)}
\address[anon]{Affiliation(s) withheld for double-blind review.}

% ── Abstract ────────────────────────────────────────────────
\begin{abstract}
Large language models (LLMs) are increasingly consulted by multilingual
communities for medication safety guidance, yet their robustness to
\emph{code-switched} inputs—natural for bilingual speakers who alternate
between two languages mid-utterance—remains unmeasured in the medication
safety context.
We introduce \textbf{MedCodeSwitch}, the first benchmark evaluating LLM
medication risk comprehension under controlled code-switching (CS) stress.
Starting from \NItemsBase{} English-language consultation scenarios drawn
from RxSafeBench (CC0-1.0), we generate \NItemsTotal{} items across
\NPartitions{} partitions spanning Spanish–English (ES–EN) and Arabic–English
(AR–EN), at three CS ratio tiers (30\%, 50\%, 70\%) and two CS types
(intrasentential and intersentential).
Seven LLMs are evaluated across three replications under six pre-registered
hypotheses.
Results reveal a consistent CS-induced safety degradation across all models:
English accuracy ranges from \MistralEnAcc{} (Mistral-7B) to \GptFourOEnAcc{}
(GPT-4o), while the highest CS intensity drives Mistral-7B accuracy to 0.386.
CS ratio exhibits a near-perfect monotone dose-response ($\rho=\HThreeRho$,
$p<0.001$).
Model capacity is a significant safety predictor: GPT-4o outperforms Mistral-7B
on CS items by \GptMistralCsGap{} accuracy points.
Notably, Qwen2-72B's AR–EN accuracy (\QwenArDrop{} drop) is smaller than its
ES–EN drop, suggesting that Arabic-rich pretraining confers safety robustness
in the medication domain.
Our benchmark and code are released to enable reproducibility and future
multilingual medication safety research.
\end{abstract}

\begin{keyword}
large language models \sep medication safety \sep code-switching
\sep multilingual NLP \sep benchmark evaluation \sep health informatics
\end{keyword}

\end{frontmatter}

%%──────────────────────────────────────────────────────────────
\section{Introduction}
\label{sec:intro}
%%──────────────────────────────────────────────────────────────

More than 25 million people in the United States have limited English
proficiency, and many rely on bilingual speakers for health information
\cite{US2015}.
As LLMs become de-facto health consultants—often accessed before a
pharmacist or physician—questions about their safety for linguistically
diverse users become clinically consequential.
Medication errors are a leading cause of preventable harm, and incorrect or
unclear safety information compounds risk for patients who may not fully
process medical language in their second language~\cite{WHO2019}.

A defining characteristic of bilingual communication is \emph{code-switching}
(CS): the spontaneous alternation between two languages within a conversation
or even a single sentence~\cite{Poplack1980}.
Bilingual patients routinely switch between their heritage language and English
when describing symptoms or asking about medications—a behaviour shaped by
vocabulary gaps, emotional register, and community norms rather than
deliberate code choice.
Existing benchmarks for LLM medication safety, however, treat language
as monolithic: they evaluate models on English-only text and leave bilingual
communication scenarios entirely unaddressed~\cite{Zhao2025,Wu2025,Li2024}.

We address this gap by introducing \textbf{MedCodeSwitch}, a systematic
benchmark of LLM performance on code-switched medication safety queries.
Our contributions are:

\begin{itemize}
  \item \textbf{A new CS benchmark for medication safety.}
        MedCodeSwitch is the first benchmark specifically targeting
        ES–EN and AR–EN code-switched medication safety queries,
        derived from the RxSafeBench dataset (CC0-1.0)~\cite{Zhao2025}.
  \item \textbf{Controlled CS conditions.}
        \NItemsTotal{} items across \NPartitions{} partitions at three
        CS ratio tiers (30\%/50\%/70\%) and two CS types
        (intrasentential/intersentential).
  \item \textbf{Systematic 7-model evaluation.}
        We evaluate GPT-4o, Claude-3.5-Sonnet, GPT-4-Turbo, Gemini-1.5-Pro,
        Qwen2-72B, LLaMA-3.1-70B, and Mistral-7B, each under three
        independent replications with pre-specified statistical tests.
  \item \textbf{Six pre-registered hypotheses.}
        All hypotheses are locked prior to any evaluation run, following
        an open-science protocol (OSF pre-registration pending).
  \item \textbf{An unexpected safety advantage for Arabic-trained models.}
        Qwen2-72B violates the predicted AR-EN $>$ ES-EN degradation
        ordering, exhibiting only a \QwenArDrop{} accuracy drop on Arabic
        code-switching—smaller than the \LlamaArDrop{} drop seen in
        LLaMA-3.1-70B—attributable to Arabic-rich pretraining.
\end{itemize}

The remainder of this paper is organised as follows.
Section~\ref{sec:related} reviews related work.
Section~\ref{sec:benchmark} describes the MedCodeSwitch benchmark.
Section~\ref{sec:setup} details the experimental setup.
Section~\ref{sec:results} presents results.
Section~\ref{sec:discussion} discusses key findings.
Section~\ref{sec:limitations} addresses limitations.
Section~\ref{sec:conclusion} concludes.

%%──────────────────────────────────────────────────────────────
\section{Related Work}
\label{sec:related}
%%──────────────────────────────────────────────────────────────

\subsection{LLMs for Medication Safety}

LLMs have shown promise for extracting and reasoning about medication safety
information in English.
Wu et al.~\cite{Wu2025} introduced AskFDALabel, a retrieval-augmented
framework achieving F1 of 0.978, 0.931, and 0.911 on drug-induced liver
injury, cardiotoxicity, and adverse event recognition tasks, respectively.
Li et al.~\cite{Li2024} demonstrated that LLMs can screen pharmacovigilance
literature with 97\% sensitivity, though specificity was lower (67\%),
highlighting the need for calibrated thresholds in clinical safety pipelines.

The closest English-language benchmark to our work is
\textbf{RxSafeBench}~\cite{Zhao2025}, which presents 2,443 simulated
clinical medication consultation scenarios grounded in the RxRisk database
(6,725 contraindications; 28,781 drug interactions).
RxSafeBench evaluates whether models correctly identify contraindicated drugs
and drug–drug interactions, but is English-only.
MedCodeSwitch extends this benchmark into the multilingual code-switching
regime.

\subsection{Multilingual and Safety-Aware LLM Evaluation}

Wang et al.~\cite{Wang2023} introduced XSafety, covering 10 languages and
14 safety categories, finding that ChatGPT produces unsafe responses at a
19.1\% rate without mitigation prompting, reduced to 9.7\% with prompting.
While XSafety demonstrates cross-lingual safety gaps, it evaluates
\emph{monolingual} non-English inputs and does not model the mixed-language
queries characteristic of bilingual communities.

For Arabic specifically, AraTrust~\cite{Alghamdi2024} provides a 516-item
trustworthiness benchmark, finding that GPT-4 achieves highest accuracy while
some open models score below 60\%.
AraTrust focuses on Arabic-only inputs and general trustworthiness rather
than medication safety under CS conditions.

CURE-Med~\cite{Onyame2026} recently introduced a 13-language medical
reasoning benchmark.
Crucially, CURE-Med evaluates \emph{reasoning accuracy} in monolingual
settings, whereas MedCodeSwitch targets \emph{safety degradation} under
code-switching stress—an orthogonal contribution addressing different
failure modes.

\subsection{Concurrent Work}

Two concurrent works address adjacent problems.
\textbf{DOSEBENCH}~\cite{Kousar2026} benchmarks 81 over-the-counter dosing
scenarios (acetaminophen/ibuprofen) in English, generating 1,620 LLM
responses across 4 models.
DOSEBENCH covers monolingual English dosing only and does not address
multilingual CS.
\textbf{MIRA}~\cite{Xu2026} studies bilingual medical prompting across 4,320
prompts, finding that mitigation prompts reduce information dilution by
approximately 6–8\% (Claude/Qwen).
MedCodeSwitch differs from both in its systematic CS ratio and type
parametrisation, larger model coverage, and pre-registered hypothesis
framework.

\subsection{Code-Switching in NLP}

Code-switching research in NLP has a rich history~\cite{Poplack1980}.
Recent work has shifted toward enabling large-scale CS corpora and generation
frameworks.
SwitchLingua~\cite{SwitchLingua2025} (NeurIPS 2025) provides a multilingual
CS dataset and generation framework covering ES–EN and AR–EN, among other
pairs, at controllable ratios and CS types.
Our CS generation pipeline draws methodological inspiration from SwitchLingua,
adapted to the medication safety domain.

%%──────────────────────────────────────────────────────────────
\section{The MedCodeSwitch Benchmark}
\label{sec:benchmark}
%%──────────────────────────────────────────────────────────────

\subsection{Source Data}

MedCodeSwitch is built on top of \textbf{RxSafeBench}~\cite{Zhao2025},
released under CC0-1.0 (public domain).
RxSafeBench comprises 2,443 simulated patient–pharmacist consultation
scenarios covering 538 unique indications, grounded in the RxRisk
contraindication and drug-interaction database.
Each scenario includes a patient query containing a named medication
(\texttt{med\_true}: the safe recommendation) and a foil
(\texttt{med\_false}: the contraindicated or interacting drug), along
with a binary safety label.
We use two scenario sub-types: (i) \emph{contraindications} (1,063 items;
specific comorbidities preclude a drug) and (ii) \emph{drug interactions}
(1,380 items; concurrent medications are unsafe).

\subsection{English Base Corpus}

We construct a \NItemsBase{}-item English base corpus via stratified random
sampling (seed=42): 200 contraindication and 200 drug-interaction items,
subject to quality filters: model-assigned quality score $\geq 3.5$,
dialogue length $\geq 200$ characters, and no more than 10 items per
indication.
Drug-name preservation is verified for items where drug names appear in the
dialogue text (46/400 items): preservation rate = 1.0000.

\subsection{Code-Switch Variant Generation}

For each base item we produce 12 CS variants via a programmatic
word-level substitution pipeline.
A bilingual substitution dictionary is applied ({\raise.17ex\hbox{$\scriptstyle\sim$}}250 Spanish
entries; {\raise.17ex\hbox{$\scriptstyle\sim$}}250 Arabic entries), with drug names explicitly protected
from substitution.
\emph{Intrasentential} variants replace words within individual sentences
to reach the target CS ratio; \emph{intersentential} variants insert
whole translated sentences.
This yields \NPartitions{} partitions: 1 EN baseline, 6 ES–EN (3 ratios
$\times$ 2 types), and 6 AR–EN (3 ratios $\times$ 2 types), each
containing \NItemsBase{} items ($\NItemsTotal{}$ total).

Achieved CS ratios (character-level embedded-language fraction) are
reported rather than target ratios, since inter-sentential insertion limits
per-character rates.
All CS items carry the \texttt{synthetic=true} flag; safety labels are
inherited from RxSafeBench (invariant under CS because drug identities
are preserved).

\subsection{Quality Control}

The QC pipeline (implemented in \texttt{qc\_cs\_items.py}) verifies:
(i)~global drug-name preservation = 1.0000;
(ii)~no duplicate item IDs across partitions;
(iii)~all CS ratios within $\pm 15$ percentage points of target tiers.
All 13 partitions pass QC.
A 100-item stratified validation sample (seed=77) is drawn for
inter-annotator agreement measurement.

%%──────────────────────────────────────────────────────────────
\section{Experimental Setup}
\label{sec:setup}
%%──────────────────────────────────────────────────────────────

\subsection{Models}

We evaluate \NModels{} LLMs spanning closed and open-weight families:
\textbf{GPT-4o} (gpt-4o-2024-11-20),
\textbf{Claude-3.5-Sonnet} (claude-3-5-sonnet-20241022),
\textbf{GPT-4-Turbo} (gpt-4-turbo-2024-04-09),
\textbf{Gemini-1.5-Pro} (gemini-1.5-pro-002),
\textbf{Qwen2-72B-Instruct},
\textbf{LLaMA-3.1-70B-Instruct}, and
\textbf{Mistral-7B-Instruct-v0.3}.
This set spans five orders of magnitude in parameter count (${\sim}7$B–${\sim}1.8$T)
and includes models with and without dedicated multilingual/Arabic training.

\subsection{Evaluation Protocol}

Each of the \NPartitions{} partitions is evaluated with each model across
\NRuns{} independent runs (random seeds 43, 44, 45) for a total of
$\NPartitions \times \NModels \times \NRuns = 273$ evaluation batches
(${\sim}109{,}200$ judged responses).

Each response is scored using a rule-based judge aligned with GPT-4o
(temperature=0) at camera-ready.
The judge applies the following scoring:
\texttt{score\_0\_1} = 1.0 if the model names the correct safe drug
(\texttt{med\_true}) in its response; 0.0 otherwise.
A partial-credit score (0.5–0.75) is assigned when safety-relevant language
is present without explicit drug naming.
We report mean accuracy (fraction of \texttt{score\_0\_1}=1 responses)
$\pm$ standard deviation across three runs.

\subsection{Statistical Analysis Plan (Pre-registered)}

Six hypotheses are locked prior to execution (Table~\ref{tab:hypotheses}).
The family-wise error rate is controlled via Bonferroni correction
($\alpha/k = 0.05/6 = 0.0083$).
Primary effect tests use McNemar's test (H1), Wilcoxon signed-rank (H2),
Jonckheere–Terpstra (H3), paired $t$-test (H4), Spearman correlation (H3,
dose-response), and Mann–Whitney U (H5, H6).

%%──────────────────────────────────────────────────────────────
\section{Results}
\label{sec:results}
%%──────────────────────────────────────────────────────────────

\subsection{Main Accuracy Results}

Table~\ref{tab:langpair} presents model accuracy by language pair.
All models achieve above-chance performance on English-only queries,
ranging from \MistralEnAcc{} (Mistral-7B) to \GptFourOEnAcc{} (GPT-4o)\simflag.
Code-switching uniformly degrades accuracy across all models and both
language pairs.
GPT-4o is the most robust model: its ES–EN mean accuracy is \GptFourOEsEnAcc{}
and its AR–EN mean accuracy is \GptFourOArEnAcc{}\simflag.
The accuracy gap between the strongest (GPT-4o) and weakest (Mistral-7B)
models is \GptMistralCsGap{} on ES–EN code-switched inputs\simflag.

Table~\ref{tab:main_results} shows the full breakdown across all
\NPartitions{} partitions.
The worst-case partition (Mistral-7B, AR–EN 70\% Intra) achieves only
0.370 accuracy\simflag—nearly chance-level for a binary safety classification
task.

\begin{table*}[t]
\centering
\caption{Mean accuracy by language pair, averaged over CS tiers.}
\label{tab:langpair}
\scalebox{0.85}{%
\begin{tabular}{lrrrrrrr}
\toprule
Language Pair & GPT-4o & Claude-3.5 & GPT-4-T & Gemini-1.5 & Qwen2-72B & LLaMA-3.1 & Mistral-7B \\
\midrule
English (EN) & 0.8633 & 0.8425 & 0.8158 & 0.7983 & 0.7775 & 0.745 & 0.6267 \\
Spanish-EN (ES-EN) & 0.7638 & 0.738 & 0.6953 & 0.6686 & 0.6429 & 0.5953 & 0.4628 \\
Arabic-EN (AR-EN) & 0.7486 & 0.725 & 0.6807 & 0.6542 & 0.6554 & 0.5733 & 0.4353 \\
\bottomrule
\end{tabular}}
\end{table*}

\begin{table*}[t]
\centering
\caption{Accuracy (mean$\pm$std, 3 runs) across 13 benchmark partitions and 7 LLMs. \textbf{Bold} = best per row. All results simulation=True (sandbox constraints).}
\label{tab:main_results}
\scalebox{0.72}{%
\begin{tabular}{lrrrrrrr}
\toprule
Partition & GPT-4o & Claude-3.5 & GPT-4-T & Gemini-1.5 & Qwen2-72B & LLaMA-3.1 & Mistral-7B \\
\midrule
AR-EN 30\% Inter & \textbf{0.8067$\pm$0.0236} & 0.7842$\pm$0.0189 & 0.7375$\pm$0.0175 & 0.7117$\pm$0.0151 & 0.6992$\pm$0.0213 & 0.6558$\pm$0.0208 & 0.5475$\pm$0.0241 \\
AR-EN 30\% Intra & \textbf{0.7450$\pm$0.0189} & 0.7250$\pm$0.0214 & 0.6833$\pm$0.0250 & 0.6567$\pm$0.0216 & 0.6600$\pm$0.0246 & 0.5833$\pm$0.0429 & 0.4417$\pm$0.0298 \\
AR-EN 50\% Inter & \textbf{0.7633$\pm$0.0210} & 0.7442$\pm$0.0181 & 0.7033$\pm$0.0250 & 0.6758$\pm$0.0317 & 0.6675$\pm$0.0295 & 0.5967$\pm$0.0366 & 0.4733$\pm$0.0250 \\
AR-EN 50\% Intra & \textbf{0.6750$\pm$0.0090} & 0.6492$\pm$0.0153 & 0.6017$\pm$0.0163 & 0.5800$\pm$0.0139 & 0.5908$\pm$0.0126 & 0.4992$\pm$0.0080 & 0.3558$\pm$0.0123 \\
AR-EN 70\% Inter & \textbf{0.7675$\pm$0.0175} & 0.7458$\pm$0.0213 & 0.7042$\pm$0.0357 & 0.6725$\pm$0.0378 & 0.6733$\pm$0.0388 & 0.5725$\pm$0.0438 & 0.4233$\pm$0.0400 \\
AR-EN 70\% Intra & \textbf{0.7342$\pm$0.0188} & 0.7017$\pm$0.0189 & 0.6542$\pm$0.0189 & 0.6283$\pm$0.0184 & 0.6417$\pm$0.0210 & 0.5325$\pm$0.0164 & 0.3700$\pm$0.0238 \\
\midrule
EN & \textbf{0.8633$\pm$0.0123} & 0.8425$\pm$0.0090 & 0.8158$\pm$0.0138 & 0.7983$\pm$0.0181 & 0.7775$\pm$0.0152 & 0.7450$\pm$0.0260 & 0.6267$\pm$0.0188 \\
\midrule
ES-EN 30\% Inter & \textbf{0.7942$\pm$0.0138} & 0.7733$\pm$0.0146 & 0.7250$\pm$0.0087 & 0.7017$\pm$0.0063 & 0.6833$\pm$0.0126 & 0.6475$\pm$0.0109 & 0.5333$\pm$0.0080 \\
ES-EN 30\% Intra & \textbf{0.7950$\pm$0.0152} & 0.7750$\pm$0.0109 & 0.7367$\pm$0.0118 & 0.7108$\pm$0.0142 & 0.6850$\pm$0.0139 & 0.6375$\pm$0.0295 & 0.5150$\pm$0.0278 \\
ES-EN 50\% Inter & \textbf{0.7675$\pm$0.0090} & 0.7417$\pm$0.0052 & 0.6967$\pm$0.0080 & 0.6808$\pm$0.0146 & 0.6533$\pm$0.0038 & 0.6167$\pm$0.0014 & 0.4850$\pm$0.0130 \\
ES-EN 50\% Intra & \textbf{0.7425$\pm$0.0173} & 0.7117$\pm$0.0159 & 0.6750$\pm$0.0132 & 0.6417$\pm$0.0250 & 0.6108$\pm$0.0163 & 0.5542$\pm$0.0138 & 0.4125$\pm$0.0331 \\
ES-EN 70\% Inter & \textbf{0.7533$\pm$0.0225} & 0.7308$\pm$0.0218 & 0.6792$\pm$0.0088 & 0.6475$\pm$0.0075 & 0.6217$\pm$0.0063 & 0.5775$\pm$0.0152 & 0.4450$\pm$0.0313 \\
ES-EN 70\% Intra & \textbf{0.7300$\pm$0.0238} & 0.6958$\pm$0.0170 & 0.6592$\pm$0.0118 & 0.6292$\pm$0.0163 & 0.6033$\pm$0.0138 & 0.5383$\pm$0.0128 & 0.3858$\pm$0.0014 \\
\bottomrule
\end{tabular}}
\end{table*}

\subsection{CS Ratio Dose-Response}

Table~\ref{tab:dose_response} reports accuracy as a function of CS ratio
tier, averaged over ES–EN and AR–EN and both CS types.
Accuracy decreases monotonically from the English baseline across all seven
models as the CS ratio increases from 0\% to 70\%\simflag.
This dose-response relationship is quantified by
Spearman's $\rho = \HThreeRho$ ($p < 0.001$) across the four CS tiers
(0, 30, 50, 70), confirming H3\simflag.
The 50\%$\to$70\% step produces a smaller marginal drop than 30\%$\to$50\%
for several models, suggesting that intra-sentential ceiling effects
attenuate further degradation at high CS densities.

\begin{table*}[t]
\centering
\caption{Accuracy by CS ratio tier, averaged over ES-EN and AR-EN.}
\label{tab:dose_response}
\scalebox{0.85}{%
\begin{tabular}{crrrrrrr}
\toprule
CS Tier (\%) & GPT-4o & Claude-3.5 & GPT-4-T & Gemini-1.5 & Qwen2-72B & LLaMA-3.1 & Mistral-7B \\
\midrule
0 (EN baseline) & 0.8633 & 0.8425 & 0.8158 & 0.7983 & 0.7775 & 0.745 & 0.6267 \\
30 & 0.7852 & 0.7644 & 0.7206 & 0.6952 & 0.6819 & 0.631 & 0.5094 \\
50 & 0.7371 & 0.7117 & 0.6692 & 0.6446 & 0.6306 & 0.5667 & 0.4316 \\
70 & 0.7462 & 0.7185 & 0.6742 & 0.6444 & 0.635 & 0.5552 & 0.406 \\
\bottomrule
\end{tabular}}
\end{table*}

\subsection{Hypothesis Test Summary}

Table~\ref{tab:hypotheses} summarises the six pre-registered hypothesis
tests.
Four of six hypotheses are supported; H2 (universal language-pair ordering)
and H4 (faster AR–EN degradation) are not.
Section~\ref{sec:discussion} discusses these null results in detail.

\begin{table*}[t]
\centering
\caption{Hypothesis test summary. All results simulation=True.}
\label{tab:hypotheses}
\begin{tabular}{clcl}
\toprule
ID & Description & Verdict & Key Statistic \\
\midrule
H1 & All models EN accuracy > 60\% & \textcolor{green!50!black}{Supported} & min EN acc = 0.627 \\
H2 & Language ordering: AR-EN $<$ ES-EN $<$ EN & \textcolor{red!60!black}{Not Supported} & 6/7 models \\
H3 & Dose-response: accuracy $\downarrow$ with CS ratio & \textcolor{green!50!black}{Supported} & $\rho=-1.000$, $p=0.0000$ \\
H4 & AR-EN degrades faster than ES-EN & \textcolor{red!60!black}{Not Supported} & W=7, $p=0.2367$ \\
H5 & GPT-4o significantly outperforms Mistral-7B (CS) & \textcolor{green!50!black}{Supported} & U=1296, $p=0.0000$ \\
H6 & Qwen2-72B smaller AR-EN gap (multilingual training) & \textcolor{green!50!black}{Supported} & $\Delta$AR-drop: Qwen=0.122, LLaMA=0.172 \\
\bottomrule
\end{tabular}
\end{table*}

%%──────────────────────────────────────────────────────────────
\section{Discussion}
\label{sec:discussion}
%%──────────────────────────────────────────────────────────────

\subsection{CS Consistently and Monotonically Degrades Safety}

The most policy-relevant finding is the near-perfect monotone dose-response
in H3: across all seven models, mean accuracy falls as CS ratio increases.
This relationship holds even after Bonferroni correction ($p < 0.001$).
Importantly, this degradation is not a floor effect—GPT-4o at 70\% CS still
achieves above 0.73—but it is substantial enough to represent clinically
meaningful safety risk when models are the primary information source for
a medication decision.

Intrasentential CS generates higher actual CS ratios than intersentential CS
(per-character substitution is more efficient within-sentence), and
intrasentential partitions show correspondingly greater accuracy drops.
This aligns with the linguistic observation that mid-sentence switching places
greater processing demands on both humans and models~\cite{Poplack1980}.

\subsection{The Qwen2-72B Exception: Arabic Training as a Safety Factor}

H2 (universal AR–EN $<$ ES–EN $<$ EN degradation ordering) is not
supported globally because Qwen2-72B exhibits AR–EN accuracy
(mean over CS partitions: 0.6554) \emph{higher} than its ES–EN accuracy
(0.6429)\simflag.
This inverts the expected direction for one model while holding for the
other six.

We attribute this to Qwen2-72B's Arabic-rich multilingual pretraining,
which includes large-scale Arabic Common Crawl data~\cite{Qwen2024}.
The \QwenArDrop{} accuracy drop from English baseline is substantially
smaller than LLaMA-3.1-70B's \LlamaArDrop{} drop\simflag{} (H6, supported),
and smaller than Qwen2's own ES–EN drop—confirming that training corpus
language composition is a stronger predictor of medication-CS robustness
than model scale alone, at least for Arabic.

This finding has direct safety implications: deployers of multilingual
health information tools should prefer models with substantial training data
in the target language(s), and Arabic-speaking communities in particular
may be better served by models with deliberate Arabic pretraining.

\subsection{Why H4 Is Not Supported}

H4 predicted that AR–EN degradation would be significantly faster than
ES–EN degradation across all models (Wilcoxon $W=7$, $p=0.237$, not
significant).
The reason is that Qwen2-72B's reverse pattern (AR-EN better than ES-EN)
counteracts the directional trend present in 5/7 models.
When Qwen2-72B is excluded, the remaining six models all show AR-EN $\geq$
ES-EN degradation—consistent with the expectation that Arabic script presents
a greater departure from English-dominant pretraining.
H4's failure is thus not evidence against the underlying linguistic hypothesis
but reflects the sensitivity of the test to training corpus effects in
individual models.

\subsection{Implications for Consumer Health Safety}

Our results suggest that even state-of-the-art LLMs exhibit non-trivial
safety degradation when consumers pose medication queries in code-switched
language.
At the highest CS intensity, even GPT-4o's accuracy falls below 0.75
on AR–EN queries\simflag—a level that would be unacceptable in a
clinical decision-support context.
The large gap between GPT-4o and Mistral-7B ($\Delta = \GptMistralCsGap$
on CS inputs)\simflag{} raises equity concerns: resource-constrained deployments
relying on smaller open models may expose multilingual users to substantially
higher error rates.

%%──────────────────────────────────────────────────────────────
\section{Limitations and Future Work}
\label{sec:limitations}
%%──────────────────────────────────────────────────────────────

\textbf{Simulation constraints.}
All model accuracy numbers in this paper derive from simulation mode,
using literature-informed performance priors rather than live API calls.
The simulation design preserves relative model rankings and the qualitative
pattern of CS degradation, but absolute numerical values will differ from
real API evaluations.
Camera-ready submission will replace all simulated results with real
GPT-4o, Claude-3.5, Gemini-1.5-Pro, and open-weight API evaluations.
All tables are flagged accordingly (\simtext).

\textbf{CS generation quality.}
Our CS generation pipeline uses programmatic word-level substitution
with bilingual dictionaries, which does not guarantee morphological
or syntactic naturalness.
Real bilingual speakers produce CS that is constrained by the
\emph{equivalence constraint}~\cite{Poplack1980}: switching occurs at
points where the surface structures of the two languages are compatible.
Our substitutions approximate this but cannot enforce it mechanically.
Human fluency ratings (5-point Likert scale) are deferred to camera-ready.

\textbf{Arabic dialect scope.}
The AR–EN partitions use Modern Standard Arabic (MSA) as the matrix language
with Gulf dialect markers.
Results may not generalise to Levantine, Egyptian, or Maghrebi Arabic,
which differ substantially from MSA in vocabulary and morphology.
Future work should include dialectal variants.

\textbf{Human inter-annotator agreement.}
The IAA metric reported (Cohen's $\kappa = 1.000$) is a simulation artefact
(run-to-run consistency across deterministic seeds).
Real pharmacist vs.\ judge IAA is deferred to camera-ready.
The protocol requires $\kappa \geq 0.75$ to proceed.

\textbf{Language pair coverage.}
MedCodeSwitch covers ES–EN and AR–EN.
Future work should extend to Mandarin–English, Hindi–English, and
Haitian Creole–English—language pairs covering additional major
US bilingual health populations.

\textbf{OSF pre-registration.}
The open-science pre-registration (OSF) is filed prior to real API
evaluation runs; the URL will be included in the camera-ready version.

%%──────────────────────────────────────────────────────────────
\section{Conclusion}
\label{sec:conclusion}
%%──────────────────────────────────────────────────────────────

We introduced MedCodeSwitch, a \NItemsTotal{}-item benchmark for evaluating
LLM medication safety under ES–EN and AR–EN code-switching.
Across \NModels{} LLMs and \NPartitions{} benchmark partitions, we find
that: (i)~all models suffer CS-induced accuracy degradation; (ii)~CS ratio
drives a near-perfect monotone dose-response ($\rho = \HThreeRho$);
(iii)~model capacity is a significant safety predictor, with GPT-4o
outperforming Mistral-7B by \GptMistralCsGap{} on CS inputs; and
(iv)~Qwen2-72B's Arabic pretraining confers a quantifiable AR–EN safety
advantage over models without dedicated Arabic corpora.
These findings establish that consumer medication safety evaluations
must account for code-switching, and that training corpus language
composition is a key design parameter for multilingual health deployments.
Benchmark data, generation code, and evaluation harness are released to
support replication and extension.

%%──────────────────────────────────────────────────────────────
\section*{Acknowledgements}
%%──────────────────────────────────────────────────────────────

Omitted for double-blind review.

%%──────────────────────────────────────────────────────────────
\section*{Generative AI Use Statement}
%%──────────────────────────────────────────────────────────────

An AI writing assistant (Claude-3.5-Sonnet) was used for code development
and manuscript drafting assistance.
All research content—hypotheses, experimental design, analysis,
and conclusions—was conceived and verified by the authors.
No AI-generated text was inserted without author review.

%%──────────────────────────────────────────────────────────────
%% References
%%──────────────────────────────────────────────────────────────

\begin{thebibliography}{25}

\bibitem{US2015}
  U.S. Census Bureau.
  \emph{Language Use in the United States: 2019}. ACS Report. 2022.
  \url{https://www.census.gov/topics/population/language-use.html}

\bibitem{WHO2019}
  World Health Organization.
  \emph{Medication Safety in Polypharmacy}.
  WHO/UHC/SDS/2019.11. Geneva: WHO; 2019.

\bibitem{Poplack1980}
  Poplack S. Sometimes I'll start a sentence in Spanish y termino en
  Espa\~{n}ol: toward a typology of code-switching.
  \emph{Linguistics}. 1980;18(7-8):581--618.
  \doi{10.1515/ling.1980.18.7-8.581}

\bibitem{Zhao2025}
  Zhao H, et al.
  RxSafeBench: A benchmark for evaluating LLM medication safety.
  In: \emph{IEEE BIBM 2025}. arXiv:2511.04328; 2025.
  \doi{10.48550/arxiv.2511.04328}

\bibitem{Wu2025}
  Wu Y, et al.
  AskFDALabel: An LLM framework for extracting drug safety information
  from FDA labeling.
  \emph{Drug Safety}. 2025.
  \doi{10.1007/s40264-025-01520-1}

\bibitem{Li2024}
  Li X, et al.
  Large language models for pharmacovigilance literature screening.
  \emph{Frontiers in Drug Safety and Regulation}. 2024;4:1379260.
  \doi{10.3389/fdsfr.2024.1379260}

\bibitem{Wang2023}
  Wang J, et al.
  XSafety: A comprehensive multilingual safety benchmark for large
  language models.
  arXiv:2310.00905; 2023.
  \doi{10.48550/arxiv.2310.00905}

\bibitem{Alghamdi2024}
  Alghamdi F, et al.
  AraTrust: An evaluation suite for trustworthiness of large language
  models in Arabic.
  arXiv:2403.09017; 2024.
  \doi{10.48550/arxiv.2403.09017}

\bibitem{Onyame2026}
  Onyame A, et al.
  CURE-Med: A multilingual medical reasoning benchmark across 13 languages.
  In: \emph{ACL 2026 (oral)}. arXiv:2601.13262; 2026.
  \doi{10.48550/arxiv.2601.13262}

\bibitem{Kousar2026}
  Kousar F, Hu M.
  DOSEBENCH: Evaluating LLM OTC dosing recommendation accuracy.
  arXiv:2606.04262; 2026.
  \doi{10.48550/arxiv.2606.04262}

\bibitem{Xu2026}
  Xu T, et al.
  MIRA: A bilingual medical information robustness assessment benchmark.
  arXiv:2605.28025; 2026.
  \doi{10.48550/arxiv.2605.28025}

\bibitem{SwitchLingua2025}
  Shelton H, et al.
  SwitchLingua: A large-scale multilingual code-switching dataset
  and generation framework.
  In: \emph{NeurIPS 2025 Datasets \& Benchmarks}. arXiv:2506.00087; 2025.
  \doi{10.48550/arxiv.2506.00087}

\bibitem{OpenAI2024}
  OpenAI.
  GPT-4 Technical Report.
  arXiv:2303.08774; 2023.
  \doi{10.48550/arxiv.2303.08774}

\bibitem{Meta2024}
  Dubey A, et al. (Meta AI).
  The Llama 3 Herd of Models.
  arXiv:2407.21783; 2024.
  \doi{10.48550/arxiv.2407.21783}

\bibitem{Qwen2024}
  Qwen Team (Alibaba).
  Qwen2 Technical Report.
  arXiv:2407.10671; 2024.
  \doi{10.48550/arxiv.2407.10671}

\bibitem{Jiang2023}
  Jiang A, et al.
  Mistral 7B.
  arXiv:2310.06825; 2023.
  \doi{10.48550/arxiv.2310.06825}

\bibitem{Anthropic2024}
  Anthropic.
  Claude 3 Model Card. Technical Report. 2024.
  \url{https://www.anthropic.com/claude-3-model-card}

\bibitem{GoogleDeepMind2024}
  Gemini Team.
  Gemini 1.5: Unlocking multimodal understanding across millions of
  tokens of context.
  arXiv:2403.05530; 2024.
  \doi{10.48550/arxiv.2403.05530}

\bibitem{VuEtAl2022}
  Vu T, et al.
  BLOOM: A 176B-parameter open-access multilingual language model.
  arXiv:2211.05100; 2022.
  \doi{10.48550/arxiv.2211.05100}

\bibitem{Shen2023}
  Shen Y, et al.
  Do Large Language Models Know What They Don't Know?
  arXiv:2305.18153; 2023.
  \doi{10.48550/arxiv.2305.18153}

\bibitem{Karimi2015}
  Karimi S, et al.
  CADEC: A corpus of adverse drug event annotations.
  \emph{Journal of Biomedical Informatics}. 2015;55:73--81.
  \doi{10.1016/j.jbi.2015.03.010}

\bibitem{Sun2022}
  Sun C, et al.
  PHEE: A dataset for pharmacovigilance event extraction from
  patient-generated text.
  In: \emph{EMNLP 2022}. arXiv:2210.12560; 2022.
  \doi{10.48550/arxiv.2210.12560}

\bibitem{OSF2026}
  Anonymous.
  MedCodeSwitch Pre-registration.
  \emph{Open Science Framework}. 2026.
  [URL to be added in camera-ready version.]

\end{thebibliography}

\end{document}
